\documentclass[aspectratio=169,10pt]{beamer}

\usetheme{metropolis}
\usepackage[T2A]{fontenc}
\usepackage[utf8]{inputenc}
\usepackage[russian]{babel}
\usepackage{amsmath,amssymb,mathtools}
\usepackage[most]{tcolorbox}
\usepackage{xcolor}

\definecolor{bg}{HTML}{0D1117}
\definecolor{bgcard}{HTML}{161B22}
\definecolor{accent}{HTML}{00FF88}
\definecolor{accentblue}{HTML}{58A6FF}
\definecolor{accentorange}{HTML}{FF7B54}
\definecolor{textmain}{HTML}{E6EDF3}
\definecolor{textsub}{HTML}{8B949E}
\definecolor{border}{HTML}{30363D}

\setbeamercolor{background canvas}{bg=bg}
\setbeamercolor{normal text}{fg=textmain,bg=bg}
\setbeamercolor{frametitle}{fg=accent,bg=bgcard}
\setbeamercolor{progress bar}{fg=accent,bg=bgcard}
\setbeamercolor{title}{fg=accent}
\setbeamercolor{subtitle}{fg=accentblue}

\setbeamerfont{title}{family=\ttfamily,series=\bfseries,size=\Large}
\setbeamerfont{frametitle}{family=\ttfamily,series=\bfseries,size=\large}
\setbeamertemplate{navigation symbols}{}

\tcbset{
  mathbox/.style={colback=bgcard,colframe=border,coltext=textmain,boxrule=0.5pt,arc=2pt,left=4pt,right=4pt,top=2pt,bottom=2pt},
  thmbox/.style={colback=bgcard,colframe=accentorange,coltext=textmain,boxrule=0pt,leftrule=2.5pt,arc=0pt,left=5pt,right=4pt,top=3pt,bottom=3pt},
  ideabox/.style={colback=bgcard,colframe=accent,coltext=textmain,boxrule=0pt,leftrule=2.5pt,arc=0pt,left=5pt,right=4pt,top=3pt,bottom=3pt}
}

\title{Название презентации}
\subtitle{Короткий подзаголовок}
\author{Имя Фамилия}
\institute{Lab or University}
\date{}

\newcommand{\topicframe}[2]{%
  \begin{frame}[plain]
    \vfill
    \begin{center}
      {\color{accent}\texttt{\small // #1}}\\[0.25cm]
      {\color{border}\rule{0.72\textwidth}{0.4pt}}\\[0.4cm]
      {\color{accent}\Huge\ttfamily #2}\\[0.35cm]
      {\color{border}\rule{0.42\textwidth}{0.3pt}}
    \end{center}
    \vfill
  \end{frame}
}

\begin{document}

\begin{frame}[plain]
  \vfill
  \begin{center}
    {\color{accent}\texttt{\small // Теория оптимизации}}\\[0.25cm]
    {\color{border}\rule{0.82\textwidth}{0.4pt}}\\[0.35cm]
    {\color{accent}\usebeamerfont{title}\inserttitle}\\[0.25cm]
    {\color{accentblue}\large\insertsubtitle}\\[0.45cm]
    {\color{border}\rule{0.55\textwidth}{0.3pt}}\\[0.25cm]
    {\color{textsub}\texttt{\small\insertauthor}}\\[0.1cm]
    {\color{textsub}\texttt{\small\insertinstitute}}\\[0.35cm]
    {\color{accent}\texttt{\small\$ run --topic curvature-aware-methods}}
  \end{center}
  \vfill
\end{frame}

\topicframe{01}{Основная тема}

\begin{frame}{Ключевая идея}
\begin{columns}[T]
\column{0.56\textwidth}
Короткий связный текст. Одна мысль на слайд.

\begin{tcolorbox}[ideabox]
Новые объекты должны быть мотивированы предыдущим слайдом.
\end{tcolorbox}

\column{0.4\textwidth}
\begin{tcolorbox}[mathbox]
$$x_{t+1}=x_t-\eta P_t^{-1}g_t$$
\end{tcolorbox}
\end{columns}
\end{frame}

\begin{frame}{Теорема}
\begin{tcolorbox}[thmbox]
Если $A\preceq B$, то шаг строится из $B$ и сохраняет нужный порядок оценки.
\end{tcolorbox}
\end{frame}

\begin{frame}[plain]
  \vfill
  \begin{center}
    {\color{accent}\texttt{\small // Questions and Discussion}}\\[0.25cm]
    {\color{border}\rule{0.72\textwidth}{0.4pt}}\\[0.35cm]
    {\color{accent}\Huge\ttfamily Спасибо!}\\[0.25cm]
    {\color{accentblue}\large Вопросы?}
  \end{center}
  \vfill
\end{frame}

\end{document}
